\documentclass[11pt]{report}
%%%%%%%%%%%%
\usepackage{graphicx}
\usepackage{pdfsync}
\usepackage{amsmath}
\usepackage{nccmath}
\usepackage{subfigure}
\usepackage{float}
\usepackage{longtable}
\usepackage[utf8]{inputenc}
%\usepackage[applemac]{inputenc}
\usepackage[letterpaper]{geometry}
\usepackage[spanish]{babel}
\usepackage{graphicx}
\usepackage{amsfonts}
\usepackage{listings}

%medidas verticales
\voffset=-1in %default 1 inch
\topmargin=2.5cm
\headheight=0pt %12pt
\headsep=0pt %25pt
\textheight=22cm
\footskip=30pt %30pt

%medidas hotizontales
\hoffset=-1in %default 1 inch
\marginparwidth=0pt %35pt
\oddsidemargin=2.5cm %31pt
\marginparsep=0pt %10pt
\textwidth=16.5cm
\usepackage{dsfont}
\usepackage{fancyhdr}
\fancyhead{}
\fancyfoot{}
\renewcommand{\headrulewidth}{0.0pt}
\renewcommand{\footrulewidth}{0.2pt}
\usepackage{hyperref}

\begin{document}
\pagestyle{fancy}

%\begin{figure}[h]
\begin{minipage}{3.8cm}
%The curret path is /home/rich 
\includegraphics[width=2.2cm]{unam.pdf} 
%Logo a color
%\includegraphics[width=2.8cm]{/home/rich/Pictures/Logos/UNAM_logo_color.jpg}
\end{minipage}
%\end{figure}
\begin{minipage}{0.42\linewidth}
\begin{center}
{\large{\textbf{Universidad Nacional Autónoma de México}}\\\textbf{Temario de Estadística Computacional y Machine Learning}\\
Prof. Jimmy Hernández Morales\\
Porf. Adjunto  Rodrigo Quijón Hipolito}
\end{center}
\end{minipage}
\vspace{1cm}

\noindent El curso pretende ser una introducción a lo que actualmente se como Machine Learning, que son un conjunto de técnicas estadísticas que junto con el desarrollo de algoritmos computacionales han tenido gran auge en los últimos años.  Todos los temas serán estudiados desde el punto de vista estadístico formal, pero haciendo mucho énfasis en la implementación computacional. 
Los lenguajes principales que usaremos serán Python o Java y como auxiliar R. Las bibliotecas esenciales que usaremos son \texttt{ Matplotlib, Scipy, Numpy, Pandas, Scikit-Learn y MLlib}. Se calificará mediante tareas, proyectos y dependiendo del tamaño del grupo exposiciones.

\subsection*{Temas}

\begin{enumerate}
\item \textbf{Regresión Lineal y Logística} 
\begin{itemize}

\item[1.1] Regresión Lineal Simple 
\item[1.2] Regresión Múltiple 
\item[1.3] Selección de Modelos 
\begin{itemize}
\item[1.3.1] Criterios AIC, BIC y divergencia de Kullback-Leibler
\end{itemize}
\item[1.4] Regresión Logística   
\end{itemize}


\item  \textbf{Aprendizaje Supervisado(Clasificación)}
\begin{itemize}
\item[2.1] Clasificadores Lineales y Gausianos. 
\item[2.2] Regresión Lineal y Logística 
\item[2.3] Discriminante Lineal(LDA) 
\item[2.3] Naive Bayes
\item[2.4] Maquinas soporte vectorial
\item[2.5] Kernelización
\item[2.6] Ejemplos, caso de estudio: redes neuronales.
\end{itemize}


\item \textbf{Modelos Aditivos, Arboles y Boosting}
\begin{itemize}
\item[3.1] Modelos aditivos Generalizados.
\item[3.2] Arboles de regresión
\item[3.3] Arboles de clasificación
\item[3.4] Métodos boosting 
\item[3.5] Ajuste boosting y modelos aditivos. 
\item[3.5] Arboles boosting
\end{itemize}


\item  \textbf{Aprendizaje no Supervisado(Clustering)}
\begin{itemize}
\item[4.1] Reglas de asociación
\item[4.2] Análisis de Clúster
\item[4.3] Componentes principales.
\end{itemize}




\item \textbf{Estadística de Altas Dimensiones y  Big Data}
\begin{itemize}
\item[5.1] Algoritmos de optimización
\begin{itemize}
\item[5.1.1] Gradiente descendiente
\item[5.1.1] Coordenadas descendientes
\end{itemize}
\item[5.2] Regresión Contraída
\begin{itemize}
\item[5.2.1] Regresión Lasso
\item[5.2.1] Regresión Ridge
\end{itemize}
\item[5.3] Correlación Espuria 
\item[5.4] Endogeneidad incidental
\begin{itemize}
\item[5.4.1] GMM
\item[5.4.2] FGMM 
 \end{itemize}
 \item[5.5] Acumulación de ruido
 \item[5.6]  Introducción al framework hadoop
 

\end{itemize}

\end{enumerate}


\begin{thebibliography}{X}

\bibitem{Baz} \textsc{Hastie, T., Tibshirani, R., Friedman, J} The Elements of Statistical Learning. Data Mining, Inference, and Prediction, 2nd ed., \textit{Springer.}

\bibitem{Baz}  \textsc{G. James, D. Witten, T. Hastie and R. Tibshirani}  An Introduction to Statistical Learning, with Application in R  \textit{Springer}, 2013).

\bibitem{Baz} \textsc{Diggle, P} 
Statistical Learning with Sparsity. The lasso and generalizations. \textit{Chapman and Hall}.


\bibitem{Baz} \textsc{T. W. Anderson} An Introduction to Multivariate Statistical Analysis, \textit{John Wiley and Sons Inc}; Edición: 3rd Edition .

\bibitem{Baz} \textsc{Michie, Spiegelhalter y Taylor} y \textsc{Petrick J. Heagarty }, Machine Learning, Neural and Statistical Classification.
\bibitem{Dan} \textsc{Alan Agresti } ,
Categorical Data Analysis,
\textit{ Wiley-Interscience Publication}.

\bibitem{Dan1} \textsc{Andreas Muller, Sarah Guido},
Introduction to Machine Learning with Python, a guide for data Scientist
\textit{Springer}. 

\end{thebibliography}
\end{document}


